\documentclass{article}
\usepackage{amssymb}
\usepackage{amsmath}
\usepackage{listings}
\usepackage{xcolor}

\lstset{
    language=Python,
    basicstyle=\ttfamily\small,
    keywordstyle=\color{blue}\bfseries,
    stringstyle=\color{teal},
    commentstyle=\color{gray}\itshape,
    frame=single,
    showstringspaces=false,
    breaklines=true
}
\begin{document}
  
  \section*{Punto 1)}
    Considere el LFSR cuyo polinomio de retroalimentación es 
      \[f(x) = x^3 + x^2 + x + 1\]

    Tomando a $S_0 = (1, 0, 0)$ como estado inicial.

    \begin{table}[h]
    \centering
      \begin{tabular}{|c|c|c|}
        \hline
        \textbf{Tick $i$} & \textbf{Estado $S_i = (y_{i+2}, y_{i+1}, y_i)$} & \textbf{Salida} \\ \hline
        0        & (1, 0, 0)                              & 0               \\ \hline
        1        & (1, 1, 0)                              & 0               \\ \hline
        2        & (0, 1, 1)                              & 1               \\ \hline
        3        & (0, 0, 1)                              & 1               \\ \hline
        4        & (1, 0, 0)                              & 0               \\ \hline
        5        & (1, 1, 0)                              & 0               \\ \hline
        6        & (0, 1, 1)                              & 1               \\ \hline
        7        & (0, 0, 1)                              & 1               \\ \hline
        8        & (1, 0, 0)                              & 0               \\ \hline
        9        & (1, 1, 0)                              & 0               \\ \hline
        10       & (0, 1, 1)                              & 1               \\ \hline
        11       & (0, 0, 1)                              & 1               \\ \hline
        12       & (1, 0, 0)                              & 0               \\ \hline
        13       & (1, 1, 0)                              & 0               \\ \hline
        14       & (0, 1, 1)                              & 1               \\ \hline
      \end{tabular}
    \end{table}

    El periodo de la sucesión es 4 y este no es un polinomio irreducible en $\mathbb{F}_2[x]$,
    pues el periodo máximo es $7$.

  \section*{Punto 2)}
    ¿Cuántos pares consecutivos se necesitan para conocer los coeficientes de retroalimentación?

    Si $m$ es el periodo del LFSR, el atacante debe conocer $2m$ pares de texto plano
    y de texto cifrado, para que con los primeros $m$ pares logre reconstruir el flujo 
    de la máquina y con los $m$ restantes construya el siguiente sistema:

    \begin{equation*}
      \begin{aligned}
        y_m &= c_{n-1} y_{m-1} \oplus c_{n-2} y_{m-2} \oplus \cdots \oplus c_0 y_{m-n} \\
        &\hspace{0.15cm}\vdots \\
        y_{2m-1} &= c_{n-1} y_{2m-2} \oplus c_{n-2} y_{2m-3} \oplus \cdots \oplus c_0 y_{2m-n-1}
      \end{aligned}
    \end{equation*}

    Y como es un sistema de ecuaciones en un cuerpo finito, tiene al menos una solución.

  \section*{Punto 3)}
    Para comprobar computacionalmente que el polinomio $f(x) = x^9 + x^4 + 1$ define un LFSR con periodo máximo, ejecutamos el siguiente código en SageMath:

    \begin{lstlisting}
      F = GF(2)
      l = F(1)
      o = F(0)

      key = [l, o, o, o, l, o, o, o, o]

      fill = [l, o, o, o, o, o, o, o, o]

      per = (2**9) - 1

      sec = lfsr_sequence(key, fill, per + 9)

      s_init = sec[:9]
      siguiente_aparicion = sec[1:].index(s_init[0]) + 1 

      print(f"Longitud teorica del periodo maximo: {per}")
      if sec[per : per+9] == s_init:
          print("El estado inicial se repite exactamente despues de 511 ticks.")
    \end{lstlisting}
\end{document}
